\begin{lemma}[Height-five outside links]
Let $N$ be normal globally maximum data for a $k$-triple
$(T,e_0,\varphi)$, suppose that the chosen height of $N$ and the Scheide
height $h(G)$ are both five, and write the trunk as
$y_i=T.\mathsf{vertex}(i)$ for $0\leq i\leq4$.  Let
$\gamma\in\overline\varphi(y_0)$.  Suppose that specified edge identities
$f_1,f_2$ and vertices
$v_1,v_2\in V(T)\setminus\{y_0,y_1,y_2,y_3,y_4\}$ satisfy
\[
 \operatorname{ends}(f_1)=\{y_1,v_1\},\quad
 \operatorname{ends}(f_2)=\{y_4,v_2\},\quad
 \varphi(f_1)=\varphi(f_2)=\gamma.
\]
Then there are distinct edge identities $a,b$ with
\[
 \operatorname{ends}(a)=\operatorname{ends}(b)=\{v_1,v_2\},
 \qquad \varphi(a)=N.\mathsf{alpha},
 \qquad \varphi(b)=N.\mathsf{beta}.
\]
\end{lemma}
\begin{proof}
Write $y_i=T.\mathsf{vertex}(i)$, put
$\alpha=N.\mathsf{alpha}$ and $\beta=N.\mathsf{beta}$, and let
$q=N.\mathsf{closingEdge}$.  Exactly as encoded by the endpoint and
support fields of \noderef{NormalGlobalMaximumTriple}, the height-five
normal path is
\[
 y_1,e_2,y_2,e_3,y_3,e_4,y_4,q,y_0.                 \tag{1}
\]
Vertex and edge injectivity in \noderef{TreeSequence} retain all displayed
identities and vertices separately.  The global field
\noderef{GlobalMaximumTriple} contains the \noderef{KTriple} condition,
so $\varphi$ is proper and leaves exactly $e_0$ uncolored.  The same
alternation and endpoint-missing argument used along (1) gives
\[
 \varphi(e_2)=\varphi(e_4)=\alpha,
 \qquad \varphi(e_3)=\varphi(q)=\beta.              \tag{2}
\]
This conclusion uses only the incidence-graph membership stored in the
normal data, the definition \noderef{TwoColorIncidenceGraph}, and
properness \noderef{IsProperEdgeColoring}.

First apply \noderef{ScheideKempePathHeightBound} to $\gamma$ missing at
$y_0$ and $\beta$ missing at $y_1$.  We obtain the full simple
$\gamma,\beta$ incidence path $P_1$ from $y_0$ to $y_1$, with all
multigraph vertices in $V(T)$ and
\[
 \operatorname{length}(P_1)/2+1\leq h(G)=5.          \tag{3}
\]
The support is the complete set \noderef{TwoColorComponent}, not merely a
chosen subwalk.  Since $q$ is the $\beta$-edge at $y_0$, properness forces
the initial segment $y_0,q,y_4$.  The specified $\gamma$-edge $f_2$ then
forces the longer segment
$y_0,q,y_4,f_2,v_2$.  At the other endpoint, $\beta$ is missing at $y_1$
and $f_1$ has color $\gamma$, so $P_1$ contains
$y_1,f_1,v_1$.

These two segments contain five distinct multigraph vertices.  The trunk
vertices $y_0,y_1,y_4$ are distinct by \noderef{TreeSequence}, and each
$v_i$ is outside the trunk by hypothesis.  Moreover $v_1\ne v_2$: the
identities $f_1,f_2$ have different endpoint pairs and hence are distinct,
while if their outside endpoints coincided they would be adjacent in the
sense of \noderef{AdjacentEdges} and would violate properness because both
have color $\gamma$.  In (3), the raw length is the number of incidence
steps.  Since the path alternates vertex-nodes and edge-nodes, the quotient
plus one is exactly its number of multigraph vertices.  Thus (3) is
saturated: $P_1$ has precisely the five vertices just listed.

Follow $P_1$ from $v_2$ after the segment ending in the $\gamma$-identity
$f_2$.  The path is not yet at its endpoint $y_1$, so it has a next
edge-node.  That identity cannot also have color $\gamma$ by properness,
and \noderef{TwoColorIncidenceGraph} says it has one of the two component
colors; hence it has color $\beta$.  Five-vertex saturation forces its
other endpoint to be $v_1$, since the remaining terminal segment is
$v_1,f_1,y_1$.  Denote this individual identity by $b$.  Then
\[
 \operatorname{ends}(b)=\{v_1,v_2\},
 \qquad \varphi(b)=\beta.                            \tag{4}
\]

We now produce the $\alpha$-identity.  First observe that the normal path
(1) is the full $\alpha,\beta$ component through $y_0$.  Indeed,
\noderef{ScheideMaximalKempePath} supplies a full simple incidence path
between $y_0$ and $y_1$.  The stored path (1) lies in that component and
already joins the same ends.  At $y_0$ the color $\alpha$ is missing and
at $y_1$ the color $\beta$ is missing, so properness gives degree one at
those endpoint vertex-nodes.  Each internal vertex-node and edge-node of
(1) already has its two path neighbors, and
\noderef{TwoColorIncidenceDegreeAtMostTwo} forbids a third neighbor.
Therefore no node of (1) has an unused incidence through which the
component could continue, and the full component has exactly the vertex-
and edge-node support recorded in (1).  In particular it contains exactly
the five trunk vertices.

The colors $\alpha$ and $\beta$ are distinct by the first conclusion of
\noderef{ScheideMaximalKempePath}.  Also $\gamma\ne\beta$, because
$\beta$ is missing at $y_1$ while the incident identity $f_1$ has color
$\gamma$; this is immediate from \noderef{MissingColors}.  Finally
$\gamma\ne\alpha$: otherwise the $\alpha$-edge-node $f_1$, incident with
$y_1$, and its outside endpoint $v_1$ would belong to the full
$\alpha,\beta$ component, contrary to the support equality just proved.

Switch the full component (1) by the definition
\noderef{TwoColorComponentSwap}.  As in
\noderef{FixedPairMaximalTashkinovTree}, global maximality follows from
the order bound \noderef{TashkinovOrder}.  The condition that $\alpha$ is
unused through $y_1$ in \noderef{TreeColorUnusedThrough} is vacuous.
Consequently \noderef{ScheideMaximalKempePath} verifies all three relevant
clauses for the switched coloring $\varphi'$: it is proper, it still has
exactly $e_0$ uncolored in the sense of
\noderef{IsOneEdgeUncolored}, and the same sequence $T$ remains a
\noderef{TashkinovTree}.  The critical, index, degree, and maximum-order
clauses are unchanged, so $(T,e_0,\varphi')$ is again globally maximum.

Under this switch, $\alpha$ becomes missing at $y_1$, the identity $q$
changes from $\beta$ to $\alpha$, and $\gamma$ remains missing at $y_0$.
The two specified identities $f_1,f_2$ and their colors are unchanged,
because $\gamma$ is distinct from both switched colors and their outside
endpoints do not belong to the full trunk component.

Apply \noderef{ScheideKempePathHeightBound} to the new globally maximum
triple with $\gamma$ missing at $y_0$ and $\alpha$ missing at $y_1$.  Its
full $\gamma,\alpha$ incidence path $P_2$ contains
\[
 y_0,q,y_4,f_2,v_2,
 \qquad y_1,f_1,v_1.                                 \tag{5}
\]
These are the same five distinct multigraph vertices as before.  The bound
$\operatorname{length}(P_2)/2+1\leq5$ therefore saturates exactly as in
(3); again, the quantity bounded is the multigraph-vertex count obtained
from the incidence length.  Following $P_2$ between the two segments in
(5), alternation and properness supply an $\alpha$-colored identity $a$
with
\[
 \operatorname{ends}(a)=\{v_1,v_2\},
 \qquad \varphi'(a)=\alpha.                          \tag{6}
\]

Both endpoints of $a$ lie outside the trunk component on which the switch
was made, so the edge-node $a$ is not in that component.  The definition
\noderef{TwoColorComponentSwap} therefore gives
$\varphi(a)=\varphi'(a)=\alpha$.  Together with (4), this proves the two
required endpoint and color assertions in the original coloring.  Finally
$a\ne b$: otherwise the equalities
$\varphi(a)=\alpha$ and $\varphi(b)=\beta$ would force
$\alpha=\beta$, contrary to the distinctness already obtained.  Thus the
returned objects are genuinely two distinct edge identities, even if they
are parallel and have the same unordered endpoint pair.
\end{proof}
